\documentclass[twocolumn]{aastex631}
\begin{document}
\title{A residual reionization ionizing-photon-budget shortfall for star-forming galaxies at z\~{}8 (lowxi systematic corner)}
\author{NebulaMind Lab (autonomous pipeline)}
\affiliation{NebulaMind Lab, \url{https://lab.nebulamind.net}}
\begin{abstract}
Reionization ionizing-photon-budget reconciliation at z\~{}8: star-forming galaxies require f\_esc=0.492 (+0.496/-0.252) to close the budget (Madau-Dickinson SFRD, log xi\_ion=25.5+/-0.15, clumping C=2-5, JWST-SFRD tail), versus indirect-proxy-inferred f\_esc=0.062 (+0.108/-0.039) from LzLCS O32/beta calibrations. Median delta(required-inferred)=+0.401 dex-frac (16-84\%: +0.136 to +0.897); 95\% of the systematic MC shows a shortfall. Conclusion: a genuine SHORTFALL remains, and the sign holds under both O32 and beta calibrations. The result is bounded by the xi\_ion x clumping x proxy-calibration systematic, not by statistics. Generated autonomously from public data (jwst) via the NebulaMind Lab runner: a bounded, reproducible, descriptive study.
\end{abstract}
\section{Introduction}
Recent studies have highlighted a potential crisis in the photon budget for reionization, suggesting that star-forming galaxies may not produce enough ionizing photons to account for the observed reionization history [Muñoz2024]. This discrepancy has sparked interest in reconciling the ionizing-photon-budget using literature-anchored calculations. Previous research has explored various aspects of this problem, including the role of absorption-dominated reionization [Davies2021], excursion set reionization models [Park2022], and assessments of the galaxy ionizing photon budget at high redshifts [Duncan2015]. Building on these efforts, we aim to investigate the reionization-photon-budget using a systematic approach grounded in published literature values.
\section{Data and method}
To address this issue, our method relies on a literature-anchored budget calculation that does not utilize survey catalog data. Instead, we adopt the cosmic star formation rate density (SFRD) from Madau \& Dickinson's (2014) analytic fitting function [Madau2017]. We also use published values for xi\_ion and O32/beta f\_esc proxy calibrations from LzLCS (Chisholm+22, Flury+22; Simmonds+24). By combining these elements, we can estimate the ionizing-photon-budget required to reconcile star-forming galaxies' contributions during reionization.
\section{Result}
Our analysis reveals that at z\~{}8, star-forming galaxies must have an escape fraction of f\_esc=0.492 (+0.496/-0.252) to close the ionizing-photon-budget when considering the Madau-Dickinson SFRD, log xi\_ion=25.5±0.15, clumping factor C=2-5, and JWST-SFRD tail. However, indirect-proxy-inferred f\_esc values derived from LzLCS O32/beta calibrations yield a significantly lower estimate of 0.062 (+0.108/-0.039). This discrepancy results in a median delta(required-inferred) of +0.401 dex-frac (16-84\%: +0.136 to +0.897), with 95\% of systematic Monte Carlo simulations indicating a shortfall. Importantly, this shortfall persists under both O32 and beta calibrations.
\begin{figure}\centering\includegraphics[width=\columnwidth]{result.png}\caption{A residual reionization ionizing-photon-budget shortfall for star-forming galaxies at z\~{}8 (lowxi systematic corner)}\end{figure}
\section{Caveats}
It is essential to acknowledge the limitations of our approach. Our analysis relies on an automated, single-selection, uncalibrated measurement that may not fully capture the complexities of reionization processes. The result is bounded by systematic uncertainties in xi\_ion, clumping factor, and proxy-calibration, rather than statistical errors. Additionally, our method does not account for potential variations in ionizing photon production efficiency or escape fraction across different galaxy populations. These caveats highlight the need for further research to refine our understanding of reionization dynamics and improve the accuracy of photon budget calculations.
\section*{References}
\footnotesize
[Muoz2024] Muñoz et al. (2024). Reionization after JWST: a photon budget crisis?. bibcode 2024MNRAS.535L..37M \par [Davies2021] Davies et al. (2021). The Predicament of Absorption-dominated Reionization: Increased Demands on Ionizing Sources. bibcode 2021ApJ...918L..35D \par [Park2022] Park et al. (2022). Calibrating excursion set reionization models to approximately conserve ionizing photons. bibcode 2022MNRAS.517..192P \par [Duncan2015] Duncan et al. (2015). Powering reionization: assessing the galaxy ionizing photon budget at z < 10. bibcode 2015MNRAS.451.2030D \par [Madau2017] Madau (2017). Cosmic Reionization after Planck and before JWST: An Analytic Approach. bibcode 2017ApJ...851...50M

\end{document}
